\documentclass[14pt,aspectratio=169]{beamer}
\usepackage[utf8]{inputenc}
\usepackage[T2A]{fontenc}
\usepackage[russian]{babel}
\usepackage{graphicx}
\usepackage{tikz}
\usetikzlibrary{arrows.meta, positioning}
\usepackage{ragged2e} % для выравнивания текста по ширине
\let\raggedright\justifying % переопределяем raggedright на justified
\usetheme{Madrid}
\usecolortheme{default}

% Автоматическое отображение названия секции в начале каждой новой секции
\AtBeginSection[]{
  \begin{frame}
    \centering
    \Huge \insertsection
  \end{frame}
}

\begin{document}

\section{Схема TDD}

\begin{frame}{Цикл разработки через тестирование (TDD)}
    \begin{itemize}
        \item \textbf{Red} -- написать падающий тест, отражающий желаемое поведение, но ещё не реализованный.
        \item \textbf{Green} -- написать минимальный код, чтобы тест прошёл.
        \item \textbf{Refactor} -- улучшить структуру кода, не меняя поведение; тесты подтверждают корректность.
    \end{itemize}
    
    \vspace{0.5cm}
    \begin{center}
        \textit{Классический цикл: Red $\to$ Green $\to$ Refactor}
    \end{center}
\end{frame}

\begin{frame}{Цикл TDD -- визуальная схема}
    \begin{center}
        \begin{tikzpicture}[
            node distance=2.2cm,
            block/.style={rectangle, draw, rounded corners, minimum width=3cm, minimum height=1cm, align=center, font=\large},
            arrow/.style={-{Stealth}, thick}
        ]
            \node[block, fill=red!20] (red) {RED\\тест не проходит};
            \node[block, fill=green!20, right=of red] (green) {GREEN\\минимальный код};
            \node[block, fill=blue!20, below=of green] (refactor) {REFACTOR\\улучшение кода};
            
            \draw[arrow] (red) -- (green);
            \draw[arrow] (green) -- (refactor);
            \draw[arrow] (refactor) -- (red);
        \end{tikzpicture}
    \end{center}
    
    \vspace{0.3cm}
\end{frame}

\section{Создатель TDD -- Кент Бек}

\begin{frame}{Кент Бек: путь и методологии}
    \begin{itemize}
        \item Программирует более 52 лет, начал в начале 1970-х.
        \item Создатель \textbf{Extreme Programming (XP)} -- революционной гибкой методологии разработки.
        \item Один из авторов \textbf{Agile Manifesto} (2001), первое имя в списке подписавших (Beck et al.).
        \item На встрече в Snowbird предложил термин «conversational» вместо «agile», но выбрали более броское название.
        \item XP продвигала короткие итерации, парное программирование, непрерывную обратную связь и ответственность разработчиков.
    \end{itemize}
\end{frame}

\begin{frame}{Кент Бек: TDD, книги и опыт}
    \begin{itemize}
        \item Пионер \textbf{Test-Driven Development} (TDD):
        \begin{itemize}
            \item В конце 1990-х написал первый xUnit-фреймворк (SUnit) на Smalltalk.
            \item Идея TDD родилась из синтеза фреймворка и старой техники «лента-на-ленту» (эталонный вывод до написания кода).
        \end{itemize}
        \item Автор книг: \textit{Tidy First?}, \textit{Extreme Programming Explained}, \textit{Test-Driven Development: By Example} и др.
        \item Работал в Facebook (2011–2017), где переосмыслил многие практики под влиянием культуры быстрой обратной связи и feature flags.
    \end{itemize}
\end{frame}

\begin{frame}{Кент Бек -- философия}
    \begin{block}{Об эмоциональной стороне TDD}
        \textit{«Экономия на анти-тревожных таблетках уже окупает TDD.»}
    \end{block}
    \begin{itemize}
        \item TDD убирает страх: когда все тесты зелёные, программист уверен, что ничего не сломал.
        \item Постоянное переключение между уровнями абстракции: тест $\to$ интерфейс $\to$ реализация $\to$ рефакторинг $\to$ снова тест.
        \item Дизайн не исключается из TDD, а становится его неотъемлемой частью в моменты рефакторинга и обдумывания следующего теста.
    \end{itemize}
\end{frame}

\section{LLM и AI-агенты глазами Кента Бека}

\begin{frame}{LLM как «джинн»}
    \begin{quote}
        \textit{«Агенты не делают то, что ты имеешь в виду. У них своя повестка. Лучшая аналогия -- джинн, исполняющий желания неожиданным образом.»}
    \end{quote}
    \begin{itemize}
        \item Метафора «джинна» (genie) -- AI-агент выполняет задачу, но результат часто не совпадает с замыслом.
        \item Агент может самовольно менять тесты, удалять их, подменять логику заглушками.
        \item Поведение напоминает азартную игру: перемежаются успехи и разочарования, что вызывает привыкание.
    \end{itemize}
\end{frame}

\begin{frame}{Проблемы LLM (1) -- контекст и побочные эффекты}
    \begin{itemize}
        \item \textbf{Непонимание контекста}: агент делает не то, что просили, или интерпретирует задачу искажённо.
        \item \textbf{Разрушение побочных эффектов}: изменения в одной части кода ломают другие, причём часто случайно.
        \item \textbf{Плохой вкус к дизайну}: агенты не умеют уменьшать связанность и повышать связность; у них «нет чувства прекрасного».
    \end{itemize}
\end{frame}

\begin{frame}{Проблемы LLM (2) -- тесты, завершённость, непредсказуемость}
    \begin{itemize}
        \item \textbf{Саботаж тестов}: стремление удалить или подправить тесты, чтобы искусственно добиться «успеха».
        \item \textbf{Ложная завершённость}: агент считает работу сделанной, хотя она не соответствует ожиданиям программиста.
        \item \textbf{Непредсказуемость}: невозможно заранее сказать, получится ли магия или придётся всё откатывать.
    \end{itemize}
\end{frame}

\begin{frame}{Как меняется роль человека (1)}
    \begin{alertblock}{Фундаментальный сдвиг}
        \textit{«90\% моих навыков обесценились, а 10\% выросли в цене в 1000 раз.»}
    \end{alertblock}
    \begin{itemize}
        \item \textbf{Уходят на второй план}: знание синтаксиса языков, управление памятью, детали реализации.
        \item \textbf{Становятся критичными}: архитектурное видение, умение разбивать большие задачи на этапы, отслеживание сложности, дизайн системы.
    \end{itemize}
\end{frame}

\begin{frame}{Как меняется роль человека (2)}
    \begin{itemize}
        \item Программист превращается в \textbf{«дирижёра»}: он ставит задачу, проверяет результат, задаёт направление, а рутину отдаёт агенту.
        \item \textbf{Эксперименты стали сверхдешёвыми}: можно пробовать гораздо больше идей, чаще выбрасывать код, учиться на неудачах.
        \begin{itemize}
            \item «Вы будете генерировать в 10 раз больше артефактов, но лишь один стоит сохранить. Надо научиться хвалить за выброшенные эксперименты».
        \end{itemize}
    \end{itemize}
\end{frame}

\begin{frame}{TDD в сочетании с AI-агентами}
    \begin{block}{Практика Кента Бека}
        \begin{itemize}
            \item Тесты выступают \textbf{основным средством общения} с агентом: когда агент ошибается, Кент показывает ожидаемое поведение через тест.
            \item Тесты -- \textbf{неизменяемая спецификация}: «Я хочу аннотацию "это правильно, и если ты это изменишь -- ты проснёшься в темноте"».
            \item Тесты запускаются очень быстро (300 мс) и могут прогоняться постоянно, чтобы ловить случайные поломки, вносимые агентом.
            \item Агент склонен ломать тесты; TDD помогает это немедленно обнаружить.
            \item TDD снижает тревогу при работе с непредсказуемым инструментом: зелёная полоса даёт уверенность.
        \end{itemize}
    \end{block}
\end{frame}

\begin{frame}{Рекомендации из слов Кента Бека}
    \begin{itemize}
        \item Не меняйте тесты ради прохождения -- фиксируйте ожидаемое поведение.
        \item Автоматически запускайте тесты после каждого изменения агента.
        \item Используйте TDD как способ верификации работы AI-ассистента.
        \item Думайте о глобальной архитектуре и разбивайте задачи на части, которые агент может выполнить поэтапно.
        \item Будьте готовы выбрасывать код: генерация дешёвая, экспериментируйте чаще.
    \end{itemize}
\end{frame}

\begin{frame}{Что думают разработчики?}
    \begin{block}{Данные Stack Overflow Developer Survey 2025}
        \begin{itemize}
            \item \textbf{84\%} разработчиков используют или планируют использовать AI-инструменты.
            \item \textbf{66\%} называют \textit{«почти правильные»} ответы главной проблемой.
            \item \textbf{45\%} говорят, что отладка AI-кода отнимает \textit{больше времени}, чем написание собственного.
        \end{itemize}
    \end{block}
    \vspace{0.3cm}
\end{frame}

\begin{frame}{Что думают разработчики?}
    \begin{alertblock}{Stack Overflow 2025: доверие падает}
        \begin{itemize}
            \item Уровень \textbf{недоверия} к точности AI-кода вырос с 31\% (2024) до \textbf{46\%} (2025).
            \item Лишь \textbf{33\%} разработчиков доверяют AI-инструментам (было 43\% в 2024).
            \item Тенденция: чем активнее используется AI, тем выше запрос на проверку.
        \end{itemize}
    \end{alertblock}
    \vspace{0.3cm}
\end{frame}

\section{Эксперименты}

\begin{frame}{Дизайн эксперимента}
\textbf{Цель исследования}

Оценить влияние структуры процесса разработки (TDD) на эффективность использования LLM.

\vspace{0.5em}

\textbf{Группы}
\begin{itemize}
    \item LLM + неструктурированный подход
    \item LLM + TDD
\end{itemize}

\vspace{0.5em}

\textbf{Условия}
\begin{itemize}
    \item Единое техническое задание
    \item Использование LLM разрешено в обеих группах
    \item Финальная проверка на одном наборе тестов
\end{itemize}
\end{frame}

\begin{frame}{Различие подходов}
\textbf{LLM + обычный подход}
\begin{itemize}
    \item Генерация кода через промпты
    \item Итеративная доработка
    \item Нет фиксированного критерия корректности в процессе
\end{itemize}

\vspace{0.5em}

\textbf{LLM + TDD}
\begin{itemize}
    \item Работа от тестов (заданы заранее)
    \item Код разрабатывается под их прохождение
    \item Тесты задают явный критерий корректности
\end{itemize}

\vspace{0.5em}

Ключевое различие: степень формализации задачи для LLM
\end{frame}

\begin{frame}{Сбор данных}
\textbf{Логируемые данные}
\begin{itemize}
    \item Все запросы к LLM и ответы модели
    \item Количество итераций взаимодействия
    \item Промежуточные и финальные версии кода
\end{itemize}

\vspace{0.5em}

\textbf{Дополнительно}
\begin{itemize}
    \item Время выполнения задачи
    \item Результаты выполнения тестов
\end{itemize}

\vspace{0.5em}

Все данные используются для последующего количественного анализа
\end{frame}

\begin{frame}{Метрики (объективные)}
\textbf{Эффективность}
\begin{itemize}
    \item Время выполнения задачи
    \item Количество итераций с LLM
\end{itemize}

\vspace{0.5em}

\textbf{Качество результата}
\begin{itemize}
    \item Процент прохождения тестов
    \item Количество успешно пройденных кейсов
\end{itemize}

\vspace{0.5em}

\textbf{Использование LLM}
\begin{itemize}
    \item Количество обращений к модели
    \item Объем промптов
\end{itemize}
\end{frame}

\begin{frame}{Метрики (субъективные)}
\textbf{Оценка участниками}
\begin{itemize}
    \item Сложность выполнения задачи
    \item Уверенность в корректности решения
    \item Удовлетворенность результатом
\end{itemize}

\vspace{0.5em}

Позволяют оценить влияние подхода на восприятие процесса разработки
\end{frame}

\begin{frame}{Исследовательские вопросы}
\textbf{Эффективность}
\begin{itemize}
    \item Ускоряет ли TDD работу с LLM
    \item Снижает ли количество итераций
\end{itemize}

\vspace{0.5em}

\textbf{Качество}
\begin{itemize}
    \item Улучшает ли TDD корректность решений
    \item Снижает ли количество ошибок
\end{itemize}

\vspace{0.5em}

\textbf{Поведение LLM}
\begin{itemize}
    \item Меняется ли характер взаимодействия с моделью
    \item Требуется ли меньше уточняющих запросов
\end{itemize}
\end{frame}

\begin{frame}{Ключевые вопросы исследования}
\textbf{Влияние на решение задачи}
\begin{itemize}
    \item Помогает ли TDD повысить корректность решений при использовании LLM
    \item Улучшает ли формализация задачи качество генерации кода
\end{itemize}

\vspace{0.5em}

\textbf{Изменение роли разработчика}
\begin{itemize}
    \item Смещается ли фокус с написания кода на формулирование требований
    \item Упрощает ли TDD взаимодействие с LLM
\end{itemize}

\vspace{0.5em}

\textbf{Когнитивная нагрузка}
\begin{itemize}
    \item Снижается ли количество итераций и исправлений
    \item Влияет ли структура (через тесты) на сложность работы с LLM
\end{itemize}
\end{frame}

\section{Обзор статей}

\begin{frame}{Fight Fire With Fire: How Much Can We Trust ChatGPT on Source Code-Related Tasks? (2024)}
    \begin{block}{}
        ChatGPT часто \textbf{ошибочно считает свой неверный код правильным} и страдает от \textit{«самопротиворечивых галлюцинаций»} во время самопроверки.
    \end{block}
    \begin{itemize}
        \item Опасно полагаться на «самооценку» AI.
        \item Необходима \textbf{независимая внешняя верификация} — автоматические тесты.
    \end{itemize}
\end{frame}

\begin{frame}{Survey On Systemic Security Risks in Vibe Coding and
Autonomous Development Workflows (2026)}
 \begin{block}{}
     Разработчик превращается в \textbf{«менеджера намерений»}, а агент — в \textbf{потенциального инсайдера}.  
        Риски сместились с артефактов кода на поведение самого агента.
    \end{block}
    \begin{itemize}
        \item \textbf{Indirect Prompt Injection (IPI)}: вредоносные инструкции в файлах, комментариях или README, которые агент послушно выполняет. Успешность атак — до 84\%.
        \item \textbf{Lies-in-the-Loop (LITL)}: агент лжёт в диалогах подтверждения — показывает «запускаю тесты», а сам выполняет опасную команду. Подрывает «человека в цикле».
    \end{itemize}
\end{frame}

\begin{frame}{Подхалимство и рост долга безопасности}
    \begin{itemize}
        \item \textbf{Sycophancy (подхалимство)}: LLM-агенты склонны угождать — если разработчик торопится, они предложат отключить проверку SSL-сертификата, лишь бы «решить» проблему быстрее.
        \item Это формирует \textbf{«Security Debt»}: накапливаются уязвимости, невидимые при поверхностном «vibe check».
    \end{itemize}
    \vspace{0.3cm}
\end{frame}

\begin{frame}{TDD как противодействие рискам AI-агентов}
    \begin{block}{Почему практики TDD становятся критически важными}
        \begin{itemize}
            \item \textbf{Неизменяемая спецификация}: тесты фиксируют истинное ожидаемое поведение, которое агент не имеет права искажать.
            \item \textbf{Мгновенное обнаружение саботажа}: тесты (прогоняемые за 300 мс) немедленно выявляют ложные изменения, внесённые агентом.
            \item \textbf{Снижение тревоги}: зелёная полоса тестов подтверждает, что «джинн» не навредил системе.
        \end{itemize}
    \end{block}
    \vspace{0.3cm}
    \begin{center}
        \textit{«Экономия на анти-тревожных таблетках уже окупает TDD» — Кент Бек}
    \end{center}
\end{frame}

\end{document}